\section{Propuesta de Solución y Metodología}
\label{ch:metodologia}

La propuesta de solución para este proyecto consiste en el diseño y entrenamiento de una arquitectura de red neuronal profunda tipo U-Net que opera en el dominio tiempo-frecuencia (STFT) para incrustar marcas de agua digitales en señales de audio. A diferencia de los métodos algorítmicos clásicos o las redes puramente temporales, la arquitectura propuesta preserva íntegramente la fase original de la señal acústica, alterando exclusivamente las magnitudes espectrales. Esta estrategia de solución garantiza simultáneamente una alta imperceptibilidad auditiva y una robustez excepcional ante ataques destructivos severos, como la compresión con pérdida (MP3) a altas tasas.

Para materializar y evaluar esta propuesta, la metodología se divide en etapas secuenciales que permiten transformar las señales de audio en representaciones tensoriales aptas para el aprendizaje profundo. Cada etapa está diseñada para garantizar la integridad de los datos y la viabilidad del entrenamiento.

\subsection{Etapa 1: Segmentación del Audio (Preprocesamiento)}
\label{sec:etapa1_segmentacion}

Debido a que las redes neuronales profundas requieren tensores de entrada con dimensiones fijas, la primera etapa consiste en una segmentación uniforme de los archivos de audio originales.

\begin{itemize}
    \item \textbf{Procedimiento:} Se cargan los archivos de audio utilizando la librería \textit{librosa}, ampliamente utilizada en el análisis de audio computacional. Cada señal se divide en fragmentos de duración objetivo $t_{target} = 5$ segundos.
    \item \textbf{Criterio de Descarte:} Para evitar el uso de técnicas de relleno (\textit{padding}) que podrían inducir sesgos, se implementó un criterio estricto donde cualquier residuo final de la señal que no complete los 5 segundos es ignorado.
    \item \textbf{Métricas de Amplitud:} Durante este proceso se registran las amplitudes máximas y mínimas de cada canal (izquierdo y derecho) para detectar posibles recortes (\textit{clipping}) y analizar el rango dinámico, lo cual es vital para determinar la ganancia de inserción de la marca de agua.
\end{itemize}

\subsection{Etapa 2: Procesamiento por Bloques (Framing) y Solapamiento (Overlap)}
\label{sec:etapa2_framing}

Para analizar la señal bajo la suposición de que es cuasi-estacionaria, se aplica una técnica de segmentación en bloques de corta duración.

\begin{itemize}
    \item \textbf{Configuración:} Se definieron bloques de $N = 512$ muestras. A una frecuencia de muestreo de 44100 Hz, esto equivale a aproximadamente 11.6 ms, tiempo suficiente para asumir estabilidad acústica.
    \item \textbf{Solapamiento:} Se utiliza un salto de $H = 256$ muestras, lo que resulta en un solapamiento (\textit{overlap}) del 50\%.
    \item \textbf{Importancia:} El solapamiento es crucial para evitar la pérdida de información en los bordes de los bloques tras la aplicación de funciones de ventana (como Hamming o Hann) y para garantizar una reconstrucción suave de la señal (\textit{Constant Overlap-Add}, COLA) tras el proceso de marcado.
\end{itemize}

\subsection{Etapa 3: Adaptación e Implementación de la Arquitectura IDEAW}
\label{sec:etapa3_diseno_red}

En esta fase se estableció la arquitectura del modelo que realiza la inserción y extracción de la marca de agua. Como punto de partida para el autoencoder y el flujo de entrenamiento, se empleó el código base del proyecto \textit{IDEAW} \cite{li2024ideaw}, proveniente de la revisión del estado del arte. No obstante, el repositorio original fue sometido a modificaciones sustanciales por parte del autor para cumplir con los objetivos específicos de esta tesina.

Las contribuciones y modificaciones arquitectónicas implementadas sobre el modelo original son las siguientes:
\begin{itemize}
    \item \textbf{Procesamiento por Espectrogramas (STFT):} El sistema fue adaptado para operar directamente sobre representaciones de tiempo-frecuencia (espectrogramas). 
    \item \textbf{Preservación Estricta de la Fase y Modificación de Magnitud (Volumen):} El repositorio original de IDEAW presentaba un defecto crítico al modificar conjuntamente todo el espectro, lo que alteraba la fase acústica y generaba distorsiones audibles. El algoritmo fue reprogramado para desacoplar el espectrograma, de modo que la red incruste la marca de agua \textit{exclusivamente} alterando la magnitud (el componente de energía o "volumen"), mientras que la fase original se extrae, se reserva y se reinyecta intacta durante la reconstrucción del audio.
    \item \textbf{Enfoque en Bajas Frecuencias (Bajos):} Se modificó el algoritmo de inserción para que la marca de agua se incruste única y selectivamente en la región de bajas frecuencias. Esto garantiza la supervivencia de la marca frente a compresores como el MP3, los cuales descartan agresivamente las altas frecuencias.
    \item \textbf{Selección Específica de Bloques:} El modelo fue reconfigurado para procesar y marcar únicamente los bloques de audio deseados, optimizando la carga computacional y asegurando que la marca se inserte solo en regiones con suficiente energía acústica.
    \item \textbf{Soporte para Canales Mono:} El pipeline original fue reprogramado para no limitarse a configuraciones estéreo, logrando operar nativamente con canales mono, lo que universaliza el sistema para todo tipo de formatos.
\end{itemize}

El entrenamiento de esta red adaptada incluye de forma intrínseca simulaciones de ruido y compresión. De este modo, el autoencoder aprende a extraer la marca de agua incluso cuando la señal ha sido degradada o atacada maliciosamente.

\begin{figure}[H]
    \centering
    \begin{tikzpicture}[node distance=2cm]
        \node (start) [startstop] {Audio Original (.wav)};
        \node (seg) [process, below of=start] {Segmentación en clips (5s)};
        \node (dec1) [decision, below of=seg, yshift=-0.5cm] {¿Cumple 5s?};
        \node (stop) [startstop, right of=dec1, xshift=2cm] {Descarte de residuo};
        \node (framing) [process, below of=dec1, yshift=-1cm] {Framing (512) y Overlap (50\%)};
        \node (stft) [process, below of=framing] {Transformada STFT};
        \node (nn) [process, below of=stft] {Autoencoder Adaptado (IDEAW)};
        \node (output) [startstop, below of=nn] {Audio Marcado / Extracción};

        \draw [arrow] (start) -- (seg);
        \draw [arrow] (seg) -- (dec1);
        \draw [arrow] (dec1) -- node[anchor=south] {No} (stop);
        \draw [arrow] (dec1) -- node[anchor=east] {Sí} (framing);
        \draw [arrow] (framing) -- (stft);
        \draw [arrow] (stft) -- (nn);
        \draw [arrow] (nn) -- (output);
    \end{tikzpicture}
    \caption{Diagrama de flujo del proceso de preprocesamiento y marcado usando IDEAW adaptado.}
    \label{fig:flujo_proceso}
\end{figure}
